\subsection{Negative Sobolev Spaces and Duality}\label{sec:negative-sobolev-duality}

According to the Riesz representation theorem, Hilbert spaces are isometric to their dual spaces. Hence every linear functional on $H^k(\Omega)$ has a representation of the form $l(v)=(u,v)_k$. However, the inner product $(u,v)_k$ does not agree with the action of $u$ as a distribution. In fact, since test functions are generally not dense in $H^k(\Omega)$, linear functionals are not necessarily distributions; there are nonzero linear functionals which vanish on all test functions. We make the following definition.

\begin{Definition}
\label{def:negative-order-sobolev-dual-spaces}
By $H^{-k}(\Omega)$, we denote the set of all linear functionals on $H_0^k(\Omega)$. Moreover, if $M$ is $\mathbb{R}^m$ or a compact manifold of class $C^k$, $k>s$, then $H^{-s}(M)$ denotes the dual space of $H^s(M)$.
\end{Definition}

Since $\mathcal{D}(\Omega)$ is dense in $H_0^k(\Omega)$, $H^{-k}(\Omega)$ is a space of distributions. As we will see in the following examples, negative Sobolev spaces contain singular distributions.

\Needspace{4\baselineskip}
\begin{Example}
\label{ex:dirac-negative-sobolev-membership}
Suppose $k>m/2$ and $\Omega\subset\mathbb{R}^m$ has the $k$-extension property and contains the origin. Then the Dirac delta is in $H^{-k}(\Omega)$. To see this we note that the Sobolev imbedding theorem ensures that $H^k(\Omega)$ (and hence $H_0^k(\Omega)$) is continuously imbedded in $C_b(\overline\Omega)$. This ensures that the delta distribution in well defined. It is also a bounded linear functional on $H_0^k(\Omega)$ since for every $u\in H_0^k(\Omega)$\[|(\delta,u)|:=|u(\bm{0})|\leq\|u\|_{\infty}\leq C\|u\|_{H^k(\Omega)}.\]
\end{Example}

\Needspace{4\baselineskip}
\begin{Example}
\label{ex:surface-charge-negative-sobolev}
Let $S$ be a smooth, bounded surface in the interior of $\Omega\subset\mathbb{R}^3$ and let $g:S\to R$ be in $L^2(S)$. (We can think of $g$ as a distribution of surface charge on $S$.) Then the distribution generated by $g$ is in $H^{-1}(\Omega)$. To see this we use the trace theorem to note that for any smooth function $\phi$ we have\[\begin{aligned}
|(g,\phi)|&=\left|\int_Sg(\bm{x})\phi(\bm{x})\,\mathrm{d}A(\bm{x})\right|\\
&\leq\|g\|_{L^2(S)}\|\phi\|_{L^2(S)}\\
&\leq\|g\|_{L^2(S)}\|\phi\|_{H^{1/2}(S)}\\
&\leq\|g\|_{L^2(S)}\|\phi\|_{H^1(\Omega)}
\end{aligned}\]Thus, the surface distribution $g$ defines a bounded linear functional on functions in $H^1(\Omega)$.
\end{Example}

We can also characterize functions in negative Sobolev spaces as derivatives of functions in positive Sobolev spaces. Let $f\in H^{-k}(\Omega)$. By the Riesz representation theorem, there is then a unique $u\in H_0^k(\Omega)$ with the property that
\begin{equation}\label{eq:negative-sobolev-riesz-representation}
(f,v)=(u,v)_k
\end{equation}
for every $v\in H_0^k(\Omega)$. How is $u$ related to $f$? From \eqref{eq:negative-sobolev-riesz-representation} we find that, for any test function $\phi$,
\begin{equation}\label{eq:negative-sobolev-riesz-derivative-pairings}
(f,\phi)=\sum_{|\alpha|\leq k}(D^\alpha u,D^\alpha\phi)=\sum_{|\alpha|\leq k}(-1)^{|\alpha|}(D^{2\alpha}u,\phi),
\end{equation}
i.e.,
\begin{equation}\label{eq:negative-sobolev-elliptic-derivative-representation}
f=\sum_{|\alpha|\leq k}(-1)^{|\alpha|}D^{2\alpha}u.
\end{equation}


For any given $f\in H^{-k}(\Omega)$, there is therefore a unique $u\in H_0^k(\Omega)$ satisfying the partial differential equation \eqref{eq:negative-sobolev-elliptic-derivative-representation}. Recall that the condition $u\in H_0^k(\Omega)$ can be interpreted as a boundary condition: $u=\partial u/\partial n=\cdots=\partial^{k-1}u/\partial n^{k-1}=0$ on $\partial\Omega$ (Theorem~\ref{thm:zero-boundary-sobolev-characterization}). Considerations similar to the one just given form the starting point of the modern existence theory for elliptic boundary-value problems; we shall return to this in Chapter~\ref{chap:linear-elliptic-equations}.

We conclude with a simple statement about differentiation of distributions in negative Sobolev spaces.

\begin{Lemma}
\label{lem:integer-order-sobolev-derivative}
Let $u\in H^k(\Omega)$, $k\in\mathbb{Z}$. Then $\partial u/\partial x_i\in H^{k-1}(\Omega)$.
\end{Lemma}

The proof follows trivially from the definitions. Lemma~\ref{lem:integer-order-sobolev-derivative} has a converse.

\begin{Lemma}
\label{lem:negative-sobolev-l-two-derivative-representation}
Let $f\in H^{-k}(\Omega)$, $k\in\mathbb{N}$. Then there exist functions $g_\alpha\in L^2(\Omega)$ such that $f=\sum_{|\alpha|\leq k}D^\alpha g_\alpha$.
\end{Lemma}

For the proof, we simply set $g_\alpha=(-1)^{|\alpha|}D^\alpha u$ in \eqref{eq:negative-sobolev-elliptic-derivative-representation}.
